\printedpage{89}
% head: A. Grothendieck : Dérivateurs, Ch. VII, Édition des Universités de Montpellier II et Paris VII
% foot: Version 22 décembre 2001

\msp{page 143}
In $1^o$ and $3^o$, we want to prove that certain arrows ($p$ and $\bar{f}$) are both \textit{proper} and \textit{smooth} - which is really ideal! I say that each time it suffices to prove one of the two, that the other follows from it (as a ``special case'') by duality. For this, let us denote by
\[
\tau \longmapsto \tau^{-}
\]
the automorphism (\textit{not anti}-automorphism) of $\Ty$ which consists in changing the sign of all the transition arrows - it is thus, at the level of objects, the unique automorphism of the free monoid Ty, which interchanges $\Delta_+$ and $\Delta_-$. To the inverse system $\underline{t}$ in $\Ty$ corresponds the inverse system $\underline{t}^-$. That said, we have an isomorphism of categories
\[
\boxed{\Ch^{\underline{t}}_{\infty}(X^o) \ \simeq\ (\Ch^{\underline{t}^-}_{\infty}(X))^o} \tag*{$(^{35})$} % Reading resolved against the supplied PDF; see corrections.csv.
\]
\footnotetext[35]{cf. page 140 for the passage from $\underline{t}$ to $\underline{t}^o$.}
and the commutative diagram
\[
\begin{tikzcd}
\Ch^{\underline{t}}_{\infty}(X^o) \arrow[r, "\sim"] \arrow[d, "p^{\underline{t}}_{X^o}"] & \Ch^{\underline{t}^-}_{\infty}(X)^o \arrow[d, "(p^{\underline{t}^-}_{X})^o"] \\
X^o \times X^o \arrow[r, equal] & (X \times X)^o.
\end{tikzcd}
\]
Thus, to say that $p^{\underline{t}}_{X}$ \texttt{[rather} $p^{\underline{t}^-}_{X}$\texttt{]} is smooth, i.e. that $(p^{\underline{t}}_{X})^o$ \texttt{[rather} $(p^{\underline{t}^-}_{X})^o$\texttt{]} is proper, is equivalent to saying that $p^{\underline{t}}_{X^o}$ is proper.

\msp{page 144}
So we are left, on this subject, with only two facts to check: the smoothness of $p : \Ch_{\infty}(X) \longrightarrow X \times X$, and that of $\bar{f} : \overline{X'} \longrightarrow X$.

a) Smoothness of $p : \Ch_{\infty}(X) \longrightarrow X \times X$. We take an arrow
\[
w : (a,b) \longrightarrow (a',b'), \qquad \text{i.e.} \quad \left\{ \begin{array}{lccc} u : & a & \longrightarrow & a' \\[1ex] v : & b & \longrightarrow & b', \end{array} \right.
\]
in the base, an object $c'$ over the target, i.e.
\[
c' \in \Ch_{\infty}(X; a', b'),
\]
and we look at the category of arrows in $\Ch_{\infty}(X)$, with target $c'$, that lift $w$,
\[
\begin{tikzcd}
\underset{\text{\scriptsize variable}}{c} \arrow[r, "\varphi"] \arrow[d, dash] & \underset{\text{\scriptsize fixed}}{c'} \arrow[d, dash] \\
(a,b) \arrow[r, "w"] & (a',b'),
\end{tikzcd}
\]
\printedpage{90}
% head: A. Grothendieck : Dérivateurs, Ch. VII, Édition des Universités de Montpellier II et Paris VII
% foot: Version 22 décembre 2001
i.e. the category of pairs $(c,\varphi)$, where
\[
\left\{ \begin{array}{l}
c \in \Ch_{\infty}(X; a, b) \\[1ex]
\varphi : c \longrightarrow c' \ \ (\in \Ch_{\infty}(X; a', b')) \\[1ex]
\text{an arrow inducing } \begin{array}[t]{l} u \ \text{ on the sources} \\ v \ \text{ on the targets,} \end{array}
\end{array} \right.
\]
\[
\begin{tikzcd}
a \arrow[rrrr, dotted, dash, "c"] \arrow[d] & {} \arrow[d, dotted] & {} \arrow[d, dotted] & {} \arrow[d, dotted] & b \arrow[d] \\
a' \arrow[rrrr, dash, "c'"] & {} & {} & {} & b'.
\end{tikzcd}
\]

\msp{page 145}
\begin{flushright}
{\footnotesize 9.11.90}
\end{flushright}

For the preceding question, concerning the arrows in the diagram
\[
\begin{tikzcd}
& X \arrow[d, "i"] & \\
& \Ch_{\infty}(X) \arrow[d, "p"] \arrow[ddl, "\sigma"'] \arrow[ddr, "\beta"] & \\
& X \times X \arrow[dl, "\mathrm{pr}_1"] \arrow[dr, "\mathrm{pr}_2"'] & \\
X & & X,
\end{tikzcd}
\]
we must absolutely specify the description of $\Ch_{\infty}(X)$ \textit{in} Cat, so we must specify the inverse system $\underline{t} = (t_n)_{n \in \mathbf{N}}$. We are in fact going to take % [sic] déscription -> description, misspelling
\[
t_n = (\underbrace{\Delta_+ \Delta_-}_{\tau_0})^{2n},
\]
with the transition morphisms described elsewhere
\[
\underbrace{t_{n+p}}_{=\tau_0^{2n+2p}} \ \longrightarrow \ \underbrace{t_n}_{\tau_0^{n}}, % [sic] \tau_0^{n} -> \tau_0^{2n}, since t_n = (\Delta_+\Delta_-)^{2n} = \tau_0^{2n}
\]
\printedpage{91}
% head: A. Grothendieck : Dérivateurs, Ch. VII, Édition des Universités de Montpellier II et Paris VII
% foot: Version 22 décembre 2001
being given by $\delta_p * \mathrm{id}_{\tau_o^{2n}} * \delta_p$, where $\delta_p : \tau_0^p \longrightarrow 1$ is the degeneracy of $\tau_0^p$. Thus we can regard $\Ch_{\infty}(X)$ as being the category of bi-infinite paths in $X$ % [sic] \tau_o^{2n} -> \tau_0^{2n}, letter o printed in the index instead of 0
\[
\cdots \longrightarrow a_{-5} \longleftarrow \overbrace{a_{-4} \longrightarrow a_{-3} \xleftarrow{u_{-2}} \overbrace{a_{-2} \xrightarrow{u_{-1}} a_{-1} \xleftarrow{u_0} \overbrace{a_0}^{\Ch_0} \xrightarrow{u_1} a_1 \xleftarrow{u_2} a_2}^{\Ch_1} \xrightarrow{u_3} a_3 \longleftarrow a_4}^{\Ch_2} \longrightarrow a_5 \longleftarrow \cdots,
\]
with the $u_i$ identities for $|i|$ large enough.

The (full) subcategory
\[
\Ch^{W}_{t_n}(X) \ = \ \Ch^{W}_{(\Delta_+\Delta_-)^{2n}}
\]
is identified with the one for which we have
\begin{quote}
$u_i$ an identity for $i \notin [-2i+1, 2i]$, i.e. the only ``significant'' $u_i$ are the $4i$ arrows $u_{-2i+1}$, $u_{-2i+2}$, $\ldots$, $u_{2i}$. % [sic] i (in [-2i+1,2i], 4i, u_{-2i+1}, u_{-2i+2}, u_{2i}) -> n, the index of t_n is n
\end{quote}

But it will be more convenient to regard $\Ch_{\infty}(X)$ as the direct limit of a
\msp{page 146}
doubly indexed family of full subcategories
\[
\Ch_{n,m}(X), \qquad n, m \in \mathbf{Z},
\]
defining this category as \textit{empty} if $n > m$, and for $n \leq m$ as consisting of the objects
\[
\left\{ \begin{array}{l}
(a_i, u_i)_{i \in \mathbf{Z}} \in \Ch^{W}_{n,m}(X) \iff u_i \text{ an identity for } i \notin [n, m-1] \\[1ex]
\overset{\mathrm{id}}{\cdots}\, a_{n-1} \dashto{\mathrm{id}} \underbrace{a_n \dashto{u_n} a_{n+1} \cdots \dashto{} a_{m-1} \dashto{u_{m-1}} a_m}_{\text{\scriptsize ``useful part'' of the path}} \dashto{\mathrm{id}} a_{m+1} \overset{\mathrm{id}}{\cdots}\,.
\end{array} \right.
\]
Thus for $n = m$, we find the full subcategory isomorphic to $X$, consisting of the ``totally degenerate'' infinite paths, i.e. such that \textit{all} the $u_i$ are identities. (This subcategory of $\Ch_{\infty}$ therefore does not depend on the choice of $n = m$.)
\[
\xrightarrow{\mathrm{id}} \underbrace{a_{-1}}_{=\, a} \xleftarrow{\mathrm{id}} \underbrace{a_0}_{=\, a} \xrightarrow{\mathrm{id}} \underbrace{a_1}_{=\, a} \xleftarrow{\mathrm{id}}
\]
All the vertices of such a path are equal to one and the same object of $X$, and the transition arrows between them are identical.

The structure of $\Ch^{W}_{n,m}(X) \ \simeq \ \Ch^{W}_{\tau_{n,m}}(X)$ depends only on the length $m-n$ of $\tau_{n,m}$ and on the \textit{parity} of $n$. If $n$ is even, the initial direction of $\tau_{n,m}$ is positive, and
\[
\tau_{n,m} = \underbrace{\Delta_+ \Delta_- \cdots}_{m-n \text{ factors in alternation}} = \left\{ \begin{array}{ll}
(\Delta_+ \Delta_-)^{\frac{m-n}{2}} & \text{if } m \text{ even} \\[1ex]
(\Delta_+ \Delta_-)^{\frac{m-n}{2}-1} \Delta_+ & \text{if } m \text{ odd}
\end{array} \right.
\]
\texttt{[rather} $\frac{m-n-1}{2}$ \texttt{instead of} $\frac{m-n}{2} - 1$\texttt{]}, whereas for $n$ odd it is just the reverse
\printedpage{92}
% head: A. Grothendieck : Dérivateurs, Ch. VII, Édition des Universités de Montpellier II et Paris VII
% foot: Version 22 décembre 2001
\msp{page 147}
\[
\tau_{n,m} = \underbrace{\Delta_- \Delta_+ \cdots}_{m-n \text{ alternating factors}} = \left\{ \begin{array}{ll}
(\Delta_- \Delta_+)^{\frac{m-n}{2}} & \text{if } m \text{ odd} \\[1ex]
(\Delta_- \Delta_+)^{\frac{m-n}{2}-1} \Delta_- & \text{if } m \text{ even}
\end{array} \right.
\]
\texttt{[rather} $\frac{m-n-1}{2}$ \texttt{instead of} $\frac{m-n}{2} - 1$\texttt{]}. The initial direction agrees with the parity of $n$, i.e. equals $(-1)^n$, whereas the terminal direction is opposite to the parity of $m$, i.e. equals $(-1)^m$ \texttt{[rather} $(-1)^{m+1}$\texttt{]}.

I note that we have $\tau_0 = (\tau_0)^o$, whence
\[
\underline{t} = \underline{t}^o, \qquad \text{but } \underline{t} \neq \underline{t}^-,
\]
and yet we have isomorphisms functorial in $X$
\[
\Ch_{\underline{t}}(X) \ \simeq \ \Ch_{\underline{t}^-}(X)
\]
respecting source and target, because of an isomorphism
\[
\underline{t} \ \simeq \ \underline{t}^- \qquad \text{in Pro}(\Ty).
\]
We can say e.g. that the \textit{translation functor} $T^p$ of amplitude $p$ (for $p \in \mathbf{Z}$) is defined on $\Ch^{W}_{\underline{t}}(X)$, by
\[
\left\{ \begin{array}{l}
T^p(((a_i),(u_i))_{i \in \mathbf{Z}}) = ((b_j),(v_j))_{j \in \mathbf{Z}} \\[1ex]
\text{with} \left\{ \begin{array}{lcl} b_j & = & a_{j-p} \\[1ex] v_j & = & u_{j-p} \end{array} \right.
\end{array} \right.
\]
and that
\[
\begin{aligned}
T^p &: \left\{ \begin{array}{lcl} \Ch_{t}(X) & \xrightarrow{\sim} & \Ch_{-t}(X) \\[1ex] \Ch_{-t}(X) & \xrightarrow{\sim} & \Ch_{t}(X) \end{array} \right\} \ \text{if } p \text{ odd} \\
T^p &: \left\{ \begin{array}{lcl} \Ch_{t}(X) & \xrightarrow{\sim} & \Ch_{t}(X) \\[1ex] \Ch_{-t}(X) & \xrightarrow{\sim} & \Ch_{-t}(X) \end{array} \right\} \ \text{if } p \text{ even.}
\end{aligned} % Reading resolved against the supplied PDF; see corrections.csv.
\]

\msp{page 148}
In other words, we can formally regard the sum category
\[
\Ch_{\pm}(X) = \Ch_{+}(X) \amalg \Ch_{-}(X)
\]
and the translation operation $T$ as an automorphism which interchanges the two pieces, and $T^p$ is nothing other than its $p$-th iterate, which interchanges the pieces $+$ and $-$ if $p$ odd, which stabilizes them if $p$ even.

\printedpage{93}
% head: A. Grothendieck : Dérivateurs, Ch. VII, Édition des Universités de Montpellier II et Paris VII
% foot: Version 22 décembre 2001

There is also the involutive automorphism of passage to the opposite path, which interchanges source and target, and which I define (in accordance with the canonical identity isomorphism
\[
\underline{t} \xrightarrow{\sim} \underline{t}^o \quad )
\]
as the symmetry $S$
\[
\left\{ \begin{array}{l}
S((a_i),(u_i)) = ((b_j),(u_j)) \ \ \text{with} \\[1ex]
\begin{array}{lcl} b_j & = & a_{-j} \\[1ex] u_j & = & u_{-j+1}. \end{array}
\end{array} \right.
\]
If I regard it as acting on $\Ch^{W}_{\pm}(X)$, it is related to the functor $T$ by
\[
\left\{ \begin{array}{rcl}
S^2 & = & \mathrm{id} \\[1ex]
STS^{-1} & = & T^{-1},
\end{array} \right.
\]
so that it is the semidirect group of $\mathbf{Z}/2\mathbf{Z}$ by $\mathbf{Z}$ (affine group of dimension 1 of the ring $\mathbf{Z}$) that acts on $\Ch_{\pm}$.

\msp{page 149}
We are embarrassed to find an automorphism $T^2$ of the situation $\Ch_{\infty}(X) \longrightarrow X \times X$, without it being obvious that in Hot, this functor induces the identity! I see no way to connect a path $c$ and $T^2 c$ by an arrow in $\Ch_{\infty}$, nor even a ``canonical'' path. But I tell myself that this automorphism is described by an automorphism, not of the inverse system % [sic] embarassé -> embarrassé, misspelling
\[
(\tau_{n,m})_{n,m \in \mathbf{Z}^o \times \mathbf{Z},\ n \leq m}
\]
(where $\mathbf{Z}^o$ denotes the ordered set \textit{opposite} to $\mathbf{Z}$), but of the \textit{pro-object} of $\Ty$ which it defines (which is of course ``the same'' for all the inverse systems cofinal on the left in $\Ty$). It already follows from this that $T^2$ induces the identity on the ind-object of Cat\,hot, $\underline{\underline{\mathrm{Ch}}}_{\infty}(X)$, and a fortiori on the corresponding ind-object of Hot, or of $\mathrm{Hot}^{\wedge}_{0}$, or finally for the object of $\mathrm{Hot}^{\wedge}_{0}$ itself. Likewise for the $\Ch_{\infty}(X; a, b)$. It follows, returning to $\Ch_{\infty}(X)$ itself and to the $\Ch_{\infty}(X; a, b)$, that the automorphism induced by $T^2$ is the identity in $\mathrm{Hot}^{\wedge}_{0}$. I hope that $\mathrm{Hot} \longrightarrow \mathrm{Hot}^{\wedge}_{0}$ is
\msp{page 150}
at least faithful, so that this will prove that $T^2$ induces the identity in Hot as well.

I return to questions $1^o$ to $4^o$ of a few days ago.

For $1^o$), we must prove that
\[
p : \Ch_{\infty}(X) \longrightarrow X \times X
\]
is smooth and proper and a Hot-fibration. By duality, $p$ proper will follow from $p$ smooth. But $\Ch_{\infty}(X)$ is $\varinjlim$ of the $\Ch_{n,m}(X)$, with $n \longrightarrow -\infty$, $m \longrightarrow +\infty$, and where we can restrict ourselves to $n$ and $m$ \textit{even.} Now we see that for $n$, $m$ even, i.e. $\tau_{n,m}$ with forward initial arrow, with final arrow